\documentclass[usepdftitle=false]{beamer}

\usepackage[frenchb]{babel}
\usepackage[T1]{fontenc}
\usepackage[utf8]{inputenc}
\usepackage{graphicx}
\usepackage{datetime}
\usepackage{eurosym}
\usepackage[]{url}
\usepackage[babel=true]{csquotes}
\usepackage{listings}
\hypersetup{
pdfauthor={Thibaud Duhautbout - Rémy Huet},
pdftitle={Formation Picaosft : La gestion de version avec Git},
pdfsubject={Formation niveau 1 : les bases},
pdfkeywords={git, gestion de version, VCS},
pdfproducer={Latex},
}

\beamertemplatenavigationsymbolsempty
\setbeamercolor{orangebox}{bg=orange,fg=black}

\newdateformat{nombres}{\THEDAY-\THEMONTH-\THEYEAR}

\title[Formation Git\_v1]{\today \\ Formation Picasoft : La gestion de version avec Git (niveau 1)}
\titlegraphic{\includegraphics[scale=.1]{picasoft_logo.png}}
\author[T. Duhautbout - R. Huet]{Thibaud {\sc Duhautbout} \\ Rémy {\sc Huet}}
\institute[Picasoft]{Association Picasoft}
\date\today

\usetheme{AnnArbor}
\usecolortheme{crane}

\AtBeginSection[]
{
	\begin{frame}
		\tableofcontents[currentsection, hideothersubsections]
	\end{frame}
}

\begin{document}

\begin{frame}
	\titlepage 
\end{frame}

\section{Introduction}

\begin{frame}{Pourquoi la gestion de version ?}
Pour tout le monde :
\begin{itemize}
\item sauvegarde incrémentale du travail
\item suivi des modifications et retour en arrière
\item partage des modifications avec d'autres personnes
\end{itemize}

\medskip

Pour les développeurs :
\begin{itemize}
\item centralisation des sources
\item collaboration simplifiée
\item possibilité de maintenir plusieurs versions
\item ...
\end{itemize}
\end{frame}

\begin{frame}{Les différents logiciels de version}
\includegraphics[height=2cm]{./imgs/logo_git.png}
\hfill
\includegraphics[height=2cm]{./imgs/logo_svn.png}
\hfill
\includegraphics[height=2cm]{./imgs/logo_mercurial.png}

\bigskip

\centering
Et plein d'autres !

37 systèmes recensés sur Wikipedia

(\url{https://en.wikipedia.org/wiki/Comparison_of_version_control_software})
\end{frame}

\begin{frame}{Petite histoire de git\ldots}

\begin{center}
\includegraphics[height=2cm]{./imgs/logo_git.png}
\end{center}

\begin{itemize}
\item Créé en 2005 par les développeurs du noyau Linux
\item Système de gestion de version distribué
\item Rapide
\item Possibilité de développements non-linéaires (branches)
\item Popularité grandissante chez les développeurs (GitHub, GitLab)
\end{itemize}
\end{frame}

\begin{frame}
\begin{beamercolorbox}[sep=5pt,center,rounded=true,shadow=true]{orangebox}
\Large
\textsc{On passe à la partie pratique !}
\end{beamercolorbox}
\begin{center}
\includegraphics[height=3cm]{./imgs/warning.jpg}
\end{center}

\centering
Dans la suite, on considère que git est installé pour toutes les machines !

\medskip

\centering
En cas de problème, n'hésitez pas à demander de l'aide aux gentils animateurs munis d'une pancarte \enquote{HELP} :)

\begin{block}{}
\centering
On ouvre un terminal ou Git Bash !
\end{block}
\end{frame}

\section{Concepts de base}

\subsection{Configuration et initialisation}

\begin{frame}[fragile]{Configuration de git et du repo}

\begin{block}{Création d'un nouveau répertoire de travail}
\verb+$ mkdir formation_git+ \\
\verb+$ cd formation_git+
\end{block}

\begin{block}{Configuration de l'identité de l'utilisateur}
Permet d'identifier l'auteur des mises à jour \\
\verb+$ git config --global/local user.name "<prénom nom>"+ \\
\verb+$ git config --global/local user.email "<adresse email>"+ \\
\verb+global+ : configuration au niveau du système \\
\verb+local+ : configuration au niveau du répertoire courant 
\end{block}

\begin{block}{Initialisation du repo git}
\verb+$ git init+
\end{block}

\end{frame}

\subsection{\'Etat du repo local}

\begin{frame}{git status}
	\begin{block}{Utilité}
		Permet de connaître à tout moment l'état d'un repo git :
		\begin{itemize}
			\item La branche sur laquelle on se situe
			\item La divergence avec un repo distant
			\item Les fichiers suivis ou non suivis
			\item Les fichiers modifiés depuis la dernière validation
			\item Les fichiers qui seront validés et ceux qui ne le seront pas.
		\end{itemize}
	\end{block}
	\begin{block}{En résumé}
		{\bf Le Saint Graal des commandes git !}
	\end{block}
\end{frame}

\subsection{Ajouter une version}

\begin{frame}{git add -- reset -- commit}
	Staging area etc
\end{frame}

\subsection{Voir l'historique}

\begin{frame}{git log}
\end{frame}

\begin{frame}{git diff}
\end{frame}

\section{Concepts avancés}

\subsection{Le HEAD}

\begin{frame}{Le HEAD}
\end{frame}

\subsection{Enregistrer les modifications locales}

\begin{frame}{git stash}
\end{frame}

\subsection{Changer de version}

\begin{frame}{git checkout}
\end{frame}

\subsection{Annuler les modifications sur un fichier précis}

\begin{frame}{git checkout -- file}
\end{frame}

\section{Les remotes}

\subsection{Principe et application avec Gitlab}

\begin{frame}{Gitlab}
	Création d'un repo
\end{frame}

\subsection{Récupérer les ajouts distants}

\begin{frame}{git clone -- git pull}
	cloner le repo de la présentation (en HTTPS)
\end{frame}

\subsection{Envoyer des modifications}

\begin{frame}{git push}
	Sur repo perso
\end{frame}

\end{document}
